% =====================================================================
% D2-NOTE — Newton-polygon universality of the Borel-singularity radius
% for polynomial continued fractions
%
% Standalone SIARC short note (Zenodo cite-target).
% Consolidates the xi_0 = d / beta_d^{1/d} cross-degree result
% currently scattered across PCF-1 v1.3, Channel Theory v1.3, and
% PCF2-SESSION-Q1.
% =====================================================================

\documentclass[11pt,reqno]{amsart}

\usepackage[margin=1in]{geometry}
\usepackage{amsmath,amssymb,amsthm,mathrsfs}
\usepackage{enumitem}
\usepackage{xcolor}
\usepackage{booktabs}
\usepackage{hyperref}
\usepackage[capitalise,noabbrev]{cleveref}
\usepackage{orcidlink}

\definecolor{darklink}{rgb}{0.0,0.2,0.5}
\hypersetup{
  colorlinks=true,
  linkcolor=darklink,
  citecolor=darklink,
  urlcolor=darklink,
  pdftitle={Newton-polygon universality of the Borel-singularity radius
            for polynomial continued fractions},
  pdfauthor={Mauricio Echizen Kubo},
  pdfkeywords={Newton polygon, Borel summability, polynomial continued
               fractions, irregular singular point, universality}
}

\theoremstyle{plain}
\newtheorem{theorem}{Theorem}[section]
\newtheorem{lemma}[theorem]{Lemma}
\newtheorem{proposition}[theorem]{Proposition}
\newtheorem{corollary}[theorem]{Corollary}
\newtheorem{conjecture}[theorem]{Conjecture}

\theoremstyle{definition}
\newtheorem{definition}[theorem]{Definition}
\newtheorem{example}[theorem]{Example}
\newtheorem{problem}[theorem]{Open problem}

\theoremstyle{remark}
\newtheorem{remark}[theorem]{Remark}

\newcommand{\PCF}{\mathrm{PCF}}
\newcommand{\dps}{\mathrm{dps}}
\newcommand{\Borel}{\mathcal{B}}

\def\version{1.0}

\title[Newton-polygon universality of the Borel-singularity radius]
      {Newton-polygon universality of the Borel-singularity radius
       for polynomial continued fractions}

\author{Mauricio Echizen Kubo}
\address{Independent researcher, Yokohama, Japan}
\email{shkubo@protonmail.com}

\date{2026-05-02 (v\version)}

\begin{document}

\begin{abstract}
We isolate, as a single citable artefact, the cross-degree
identity
\[
  \xi_0(b) \;=\; \frac{d}{\beta_d^{1/d}}
\]
for the leading Borel-plane singular distance of the formal
generating function $f(z) = \sum_{n\ge 0} Q_n z^n$ of a
polynomial continued fraction $(1, b)$ of degree $d$ with
leading coefficient $\beta_d > 0$. The identity is currently
distributed across three SIARC artefacts: the canonical
$d = 2$ Newton-polygon proof (Channel Theory v1.3,
\cite{siarc_channel_theory_v13}), the $d = 4$
verification at $\dps = 80$ across $8$ quartic
representatives with spread $0$
(PCF2-SESSION-Q1, \cite{siarc_pcf2_session_q1}),
and the original $d = 2$ baseline (PCF-1 v1.3,
\cite{siarc_pcf1_v13}). We frame the resulting picture
in three explicit epistemic tiers:
(i) \emph{proven} at $d = 2$ by the Newton-polygon argument
recalled here;
(ii) \emph{verified} (in the empirical sense: a measurement
at $\dps = 80$ with spread $0$ across $8$ representatives) at
$d = 4$;
(iii) \emph{conjectured} for general $d \ge 2$
(Conjecture~3.3.A$^{*}$ of \cite{siarc_channel_theory_v13}),
with the $d = 3$ direct Newton-polygon test deferred to the
operator queue. The note is purely consolidatory; no new
numerical pipeline is run.
\end{abstract}

\maketitle

\section{Setup}
\label{sec:setup}

Throughout this note $(1, b)$ denotes a polynomial continued
fraction with constant numerator $a_n \equiv 1$ and
polynomial denominator
\[
  b(n) \;=\; \beta_d\,n^d + \beta_{d-1}\,n^{d-1} + \dots
              + \beta_1\,n + \beta_0,
  \qquad \beta_d > 0,
\]
with $d \ge 2$. Let $Q_n$ denote the partial denominators
of the PCF, satisfying the three-term recurrence
$Q_n = b(n)\,Q_{n-1} + Q_{n-2}$ with $Q_{-1} = 0$,
$Q_0 = 1$. The formal generating function
\[
  f(z) \;=\; \sum_{n\ge 0} Q_n\,z^n
\]
satisfies the order-$2$ inhomogeneous linear ODE
\begin{equation}
\label{eq:Lf}
  L_d f \;:=\;
   \bigl(\,1 \,-\, z\,B_d(\theta + 1)
            \,-\, z^2\,\bigr)\,f \;=\; 1,
  \qquad \theta := z\,\partial_z,
\end{equation}
with $B_d$ the $d$-th degree polynomial encoding $b$ on the
shift-rescaled Euler operator $\theta$. Equation
\eqref{eq:Lf} has a regular point at $z = \infty$ and an
irregular singular point at $z = 0$; the Borel-plane
singular distance $\xi_0$ is the modulus of the leading
characteristic root at $z = 0$ along the slope-$1/d$ edge
of the Newton polygon of the homogeneous part of
\eqref{eq:Lf}. The scope of this note is fixed by
$\beta_d > 0$ and $d \ge 2$; the $\beta_d < 0$ branch and
the so-called ``$d = 2$ anomaly'' Galois bin
(\cite[Remark~5.E]{siarc_pcf1_v13}) are out of scope and
are recorded as open in \Cref{sec:open}.

\section{The proven case \texorpdfstring{$d = 2$}{d=2}}
\label{sec:d2}

The $d = 2$ Newton-polygon argument is in canonical form in
\cite[Proposition~3.3.A and Worked Example
\S3.3]{siarc_channel_theory_v13}, with the same content
recorded earlier in
\cite[Theorem~5, \S5]{siarc_pcf1_v13}. We follow the
Channel Theory v1.3 idiom because it carries the slope and
characteristic-polynomial framing needed for the cross-degree
extension in \Cref{sec:d4,sec:conjecture}.

\begin{proposition}[Universality of $\xi_0$, degree $2$;
{\cite[Prop.~3.3.A]{siarc_channel_theory_v13}}]
\label{prop:xi0-d2}
For every degree-$2$ PCF $(1, b)$ in scope with
$b(n) = \beta_2 n^2 + \beta_1 n + \beta_0$ and
$\beta_2 > 0$, the leading Borel-plane singularity of the
formal generating function $f(z) = \sum_{n \ge 0} Q_n z^n$
at $z = 0$ is located at distance
\[
  \xi_0 \;=\; \frac{2}{\sqrt{\beta_2}}
\]
from the origin, and the secondary indicial exponent of the
formal solution at the irregular singular point
$u = \sqrt{z} = 0$ is
\[
  \rho \;=\; -\frac{3}{2} \,-\, \frac{\beta_1}{\beta_2}.
\]
\end{proposition}

\begin{proof}[Sketch (after
{\cite[\S3.3]{siarc_channel_theory_v13}})]
The Newton polygon of the homogeneous part of \eqref{eq:Lf}
at $z = 0$, with $d = 2$, is read off from the lattice
points
$\{(0,0),(1,0),(1,1),(1,2),(2,0)\}$
(weight = order of $\theta$, height = order of $z$). The
lower convex hull has a single non-trivial edge
$(0,0) \to (1,2)$ of slope $1/2$ and multiplicity $2$,
corresponding to a Gevrey-$2$-in-$z$ irregular singularity.
Substituting $z = u^2$ and
$\theta = (u/2)\,\partial_u$, the operator becomes
Gevrey-$1$-in-$u$, and the level-$1/u$ trans-series ansatz
\[
  f_\pm(u) \;=\; \exp\!\Bigl(\frac{c}{u}\Bigr)\,u^{\rho}\,
  \Bigl(\,1 + \sum_{k\ge 1} a_k u^k\,\Bigr)
\]
produces the characteristic polynomial
\[
  \chi(c) \;=\; 1 \,-\, \frac{\beta_2}{4}\,c^2
\]
at leading order, whose two roots are
$c = \pm\,2/\sqrt{\beta_2}$. The positive root pins
$\xi_0 = 2/\sqrt{\beta_2}$. The $u^1$ coefficient of the
trans-series equation pins the indicial polynomial, giving
$\rho = -3/2 - \beta_1/\beta_2$. The Birkhoff recursion for
the $a_k$ then proceeds row-by-row in the $u^k$
coefficient. The full computation, including the worked
$V_{\mathrm{quad}}$ example and the
$(\alpha_1, \alpha_0)$-independence of $\xi_0$, is
\cite[\S3.3 Worked Example]{siarc_channel_theory_v13}.
\end{proof}

\begin{remark}[$(\alpha_1, \alpha_0)$-independence at $d = 2$]
\label{rem:numerator-indep}
The numerator polynomial $a$ enters the operator $L_2$ only
at the $z^2$ Newton-polygon vertex $(2,0)$, i.e.\ at order
$u^4$ in the $u$-uniformised operator. Therefore $\xi_0$
(read off at order $u^0$) and $\rho$ (read off at order
$u^1$) are independent of the numerator. The first formal
invariant that does depend on $\alpha_1, \alpha_0$ is $a_2$,
which is the splitting witness in
\cite[Prop.~3.3.B]{siarc_channel_theory_v13} for the
QL01/QL02 family pair (identical $\xi_0 = 2$ and $\rho = -5/2$,
disagreeing $a_2$).
\end{remark}

\section{The verified case \texorpdfstring{$d = 4$}{d=4}}
\label{sec:d4}

At $d = 4$, the same Newton-polygon construction predicts a
single slope-$1/4$ edge of multiplicity $2$ with
characteristic polynomial $1 - (\beta_4/4^4)\,c^4$, hence
the prediction $\xi_0 = 4/\beta_4^{1/4}$. This prediction
has been tested computationally in
\cite{siarc_pcf2_session_q1} and is consistent with the data
described below.

\begin{proposition}[Empirical record at $d = 4$;
{\cite[Prop.~3.3.A$'$]{siarc_channel_theory_v13},
\cite{siarc_pcf2_session_q1}}]
\label{prop:xi0-d4}
For a degree-$4$ PCF $(1, b)$ with
$b(n) = \beta_4 n^4 + \mathrm{l.o.t.}$ and $\beta_4 > 0$,
the Newton polygon of \eqref{eq:Lf} at $z = 0$ has lattice
points
$\{(0,0),(1,0),(1,1),(1,2),(1,3),(1,4),(2,0)\}$, with
lower-left convex hull supporting the edge
$E:(0,0) \to (1,4)$ of slope $1/4$. The WKB ansatz
$f \sim \exp(c/u)$ with $z = u^4$ then pins the
characteristic equation
\[
  \chi(c) \;=\; 1 \,-\, \frac{\beta_4}{4^4}\,c^4 \;=\; 0,
\]
whose positive real root is $c = 4/\beta_4^{1/4}$. The
identity $\xi_0(b) = 4/\beta_4^{1/4}$ is consistent with
the PCF2-SESSION-Q1 measurement: across $8$ quartic
representatives drawn from the four trichotomy bins of
the session ($S_4$ generic, $D_4$ Eisenstein anchor,
$V_4$ biquadratic anchor, plus $\beta_4 = 7$ sample), the
inferred $c(4)$ has spread $0$ at $\dps = 80$.
\end{proposition}

The detailed Newton-polygon construction and the per-family
table are in
\cite[\texttt{newton\_polygon\_d4.tex}]{siarc_pcf2_session_q1};
the PCF2-SESSION-Q1 \texttt{claims.jsonl} entry
$\texttt{Q1-B}$ carries the SHA-$256$ output hash
\texttt{1ad90f60$\dots$d5b10bc7} of the underlying script
\texttt{session\_q1\_newton.py} and is the canonical AEAL
anchor for the $d = 4$ measurement.

\begin{remark}[Falsification of the v1.1 candidate
$c(d) = 2\sqrt{(d-1)!}$;
{\cite[Remark~3.3.E]{siarc_channel_theory_v13}}]
\label{rem:falsification}
Version~1.1 of the Channel Theory outline conjectured the
candidate $c(d) = 2\sqrt{(d-1)!}$ for the universality
constant in $\xi_0(b) = c(d) / \beta_d^{1/d}$. That
candidate matches the proven $d = 2$ value
($c(2) = 2\sqrt{1!} = 2$) but is ruled out at $d = 4$ by
the PCF2-SESSION-Q1 measurement: the $\dps = 80$ data
narrow the empirical $c(4)$ to a single value $4$ with
spread $0$ across $8$ representatives, against the v1.1
prediction $c(4) = 2\sqrt{3!} = 2\sqrt{6} \approx 4.899$
-- a $\approx 22\%$ disagreement, four orders of magnitude
above the per-representative numerical resolution. The
linear form $c(d) = d$ of \Cref{conj:xi0-univ} is the form
consistent with both the proven $d = 2$ value and the
$d = 4$ measurement.
\end{remark}

\section{The conjectured general-\texorpdfstring{$d$}{d} case}
\label{sec:conjecture}

\begin{conjecture}[Cross-degree universality of $\xi_0$;
{\cite[Conj.~3.3.A$^{*}$]{siarc_channel_theory_v13}}]
\label{conj:xi0-univ}
For every PCF $(1, b)$ in scope with $\deg b = d \ge 2$ and
$b(n) = \beta_d n^d + \mathrm{l.o.t.}$ ($\beta_d > 0$), the
leading characteristic root of the irregular singularity at
$z = 0$ of the operator $L_d$ in \eqref{eq:Lf} is
\[
  \xi_0(b) \;=\; \frac{d}{\beta_d^{1/d}}.
\]
\end{conjecture}

\noindent\emph{Status.} \emph{Proven} at $d = 2$
(\Cref{prop:xi0-d2}). \emph{Verified} at $d = 4$ at
$\dps = 80$ with spread $0$ across $8$ representatives
(\Cref{prop:xi0-d4}). Pending direct Newton-polygon test
at $d = 3$ (\textsf{op:xi0-d3-direct}; see
\Cref{sec:open}). \emph{Conjectured} for general $d \ge 2$.

\subsection*{Heuristic structural picture}

The structural Newton-polygon picture that supports the
conjecture is as follows; we record it as a heuristic, not
as a proof. For any monic polynomial $b$ of degree $d$, the
operator $L_d$ in \eqref{eq:Lf} has Newton polygon at
$z = 0$ supported on the lattice points
\[
  \{(0,0)\} \;\cup\; \{(1, k) : 0 \le k \le d\}
            \;\cup\; \{(2, 0)\},
\]
and the lower-left convex hull has a single non-trivial edge
$E:(0,0) \to (1, d)$ of slope $1/d$ and multiplicity~$2$.
Substituting $z = u^d$ with $\theta = (u/d)\,\partial_u$ and
applying the trans-series ansatz $f \sim \exp(c/u)$ pins the
characteristic polynomial along $E$ to
\[
  \chi(c) \;=\; 1 \,-\, \frac{\beta_d}{d^d}\,c^d
\]
at leading order, whose unique positive real root is
$c = d / \beta_d^{1/d}$. This is the structural reason why
$c(d) = d$ is the form consistent with both the proven
$d = 2$ value and the $d = 4$ measurement, and why the v1.1
candidate $c(d) = 2\sqrt{(d-1)!}$ is ruled out
(\Cref{rem:falsification}).

The picture above does \emph{not} discharge the general-$d$
case: the indicial-polynomial analysis fixing the secondary
exponent $\rho_d$, and the Birkhoff recursion delivering the
formal coefficients $a_k$ at $d \ge 3$, are not written out
here and are open even at $d = 3$
(see \Cref{sec:open}, item~(a)).

\section{Status of \texorpdfstring{$d = 3$}{d=3}}
\label{sec:d3}

The $d = 3$ direct Newton-polygon test is the next discrete
step after the $d = 2$ proof and the $d = 4$ measurement.
It is recorded as an operator queue entry in
\cite[\S9, \textsf{op:xi0-d3-direct}]{siarc_channel_theory_v13}
and is currently \emph{deferred}, not run. The recipe is one
Newton-polygon construction per Galois bin on a cubic
representative $b(n) = \beta_3 n^3 + \dots$, with the
prediction $\xi_0 = 3/\beta_3^{1/3}$ to be checked at
$\dps = 80$ against the leading characteristic root of the
slope-$1/3$ edge of \eqref{eq:Lf}; the current tractability
estimate is $1$--$2$ hours of computation
\cite[\S9]{siarc_channel_theory_v13} with no halt condition
expected.

A successful $d = 3$ measurement would narrow the interval
of candidates for $c(d)$ further but would not by itself
discharge the general-$d$ case
(\Cref{conj:xi0-univ}); a structural proof of the
indicial / Birkhoff layer at $d \ge 3$ remains open.

\section{Open questions}
\label{sec:open}

\begin{enumerate}[label=(\alph*),leftmargin=2.2em]

\item \textsf{op:xi0-d3-direct}
(\cite[\S9]{siarc_channel_theory_v13}): direct
Newton-polygon test of $\xi_0(b) = 3/\beta_3^{1/3}$ on at
least one cubic representative per Galois bin, at
$\dps = 80$. Predicted runtime: $1$--$2$ hours per bin. No
halt condition is expected; a measurement consistent with
the prediction would narrow the empirical interval for
$c(d)$ to $\{2, 3, 4\}$ at $d \in \{2, 3, 4\}$ and would
support \Cref{conj:xi0-univ} along the linear arithmetic
progression. A measurement inconsistent with the prediction
would falsify the conjecture and trigger the corresponding
SIARC halt condition.

\item The $d \ge 5$ case. No empirical record exists for
$d \ge 5$; the Newton-polygon construction
(\Cref{sec:conjecture}) extends in principle, but neither
the slope-$1/d$ edge analysis at finite $\beta_d$ nor the
secondary indicial exponent $\rho_d$ has been written out.
A combinatorial generalisation of \Cref{prop:xi0-d2} to all
$d \ge 2$ -- including a uniform proof of the
$(\alpha_1, \dots, \alpha_{d-2})$-independence of $\xi_0$
(\Cref{rem:numerator-indep}) -- is the natural next
target.

\item The $\beta_d < 0$ branch and the $d = 2$ anomaly
Galois bin
(\cite[Remark~5.E]{siarc_pcf1_v13}). The slope-$1/d$ edge
analysis here predicts $\xi_0 = d / |\beta_d|^{1/d}$ at the
level of moduli, but the Stokes-cocycle phase and the
character of the corresponding Borel singularity are
expected to differ. A separate analysis is open.

\end{enumerate}

\section*{AI Disclosure}

The SIARC v1.3 cohort methodology was used in drafting this
note: GitHub Copilot (powered by Anthropic Claude Opus~4.7)
and Anthropic Claude were used for prose drafting and
consistency-checking; all theorem statements, conjecture
formulations, and editorial decisions remain the author's,
and every numerical claim traces to an AEAL entry in
\texttt{claims.jsonl} with a SHA-$256$ output hash on the
SIARC bridge under \texttt{sessions/2026-05-02/D2-NOTE-DRAFT/}.

\bibliographystyle{plain}
\bibliography{annotated_bibliography}

\end{document}
